\newcommand{\eps}{\epsilon}
\newcommand{\veps}{\varepsilon}
\newcommand{\Qed}{\begin{flushright}\qed\end{flushright}}

\newcommand{\parinn}{\setlength{\parindent}{1cm}}
\newcommand{\parinf}{\setlength{\parindent}{0cm}}

% \newcommand{\norm}{\|\cdot\|}
\newcommand{\inorm}{\norm_{\infty}}
\newcommand{\opensets}{\{V_{\alpha}\}_{\alpha\in I}}
\newcommand{\oset}{V_{\alpha}}
\newcommand{\opset}[1]{V_{\alpha_{#1}}}
\newcommand{\lub}{\text{lub}}
\newcommand{\del}[2]{\frac{\partial #1}{\partial #2}}
\newcommand{\Del}[3]{\frac{\partial^{#1} #2}{\partial^{#1} #3}}
\newcommand{\deld}[2]{\dfrac{\partial #1}{\partial #2}}
\newcommand{\Deld}[3]{\dfrac{\partial^{#1} #2}{\partial^{#1} #3}}
\newcommand{\der}[2]{\frac{\mathrm{d} #1}{\mathrm{d} #2}}
% \newcommand{\ddd}[3]{\frac{\mathrm{d}^{#3} #1}{\mathrm{d}^{#3} #2}}
\newcommand{\lm}{\lambda}
\newcommand{\uin}{\mathbin{\rotatebox[origin=c]{90}{$\in$}}}
\newcommand{\usubset}{\mathbin{\rotatebox[origin=c]{90}{$\subset$}}}
\newcommand{\lt}{\left}
\newcommand{\rt}{\right}
\newcommand{\bs}[1]{\boldsymbol{#1}}
\newcommand{\exs}{\exists}
\newcommand{\st}{\strut}
\newcommand{\dps}[1]{\displaystyle{#1}}
\newcommand{\id}{\text{id}}
\newcommand{\bb}{\mathbb}

\newcommand{\sol}{\textcolor{myr}{\textbf{[RASCUNHO --- substituir]}}\par\smallskip}
\newcommand{\solve}[1]{\setlength{\parindent}{0cm}\textbf{\textit{Solution: }}\setlength{\parindent}{1cm}#1 \Qed}

\newcommand{\qslista}[3][1]{%
	\begin{question}{\hypersetup{linkcolor=white}\hyperlink{sol:lista#1:#2}{Lista #1, Exercício #2}}{}%
		\hypertarget{enun:lista#1:#2}{}%
		#3%
		\par\smallskip\hfill{\footnotesize\hypersetup{linkcolor=doc}\hyperlink{sol:lista#1:#2}{\textbf{[Ver Solução $\to$]}}}%
	\end{question}%
}

\newcommand{\currentlista}{1}
\newcounter{exq}
\newenvironment{exercicio}[1][]
  {\stepcounter{exq}%
   \hypertarget{sol:lista\currentlista:\theexq}{}%
   \begin{solucaobox}{\hypersetup{linkcolor=white}\hyperlink{enun:lista\currentlista:\theexq}{Solução do Exercício \theexq}}%
   \noindent\hfill{\footnotesize\hypersetup{linkcolor=doc}\hyperlink{enun:lista\currentlista:\theexq}{\textbf{[Enunciado no Capítulo $\leftarrow$]}}}\par\smallskip
  }
  {\end{solucaobox}}

\newenvironment{obs}[1][]
  {\par\smallskip\begin{note}\ifstrempty{#1}{}{\textbf{#1:}\par\smallskip}}
  {\end{note}\par\smallskip}

% Notação usada nas soluções (chapter/solucoes-lista1.tex e solucoes-lista2.tex)
\newcommand{\R}{\bbR}
\newcommand{\N}{\bbN}
\newcommand{\Z}{\bbZ}
\newcommand{\K}{\bbK}
\newcommand{\Ker}{\mathcal{N}}
\newcommand{\Ima}{\operatorname{Im}}
\newcommand{\dist}{\operatorname{dist}}
